\subsection{The isolated degree-thirty-six Bernoulli certificate}
\label{sec:bernoulli-residual-prime36}

\begin{proposition}[K17-BERNOULLI-8361790201]
\label{prop:bernoulli-residue36}
For $p=8361790201$ and $n=p-19=8361790182$,
\[
                         B_n\equiv1827838999\ne0\pmod p.
 \tag{BC36.1}
\]
\end{proposition}

\begin{proof}
Use the already proved finite congruence (BC.3) with multiplier $g=2$.
It does not require this multiplier to be a primitive root; only
$2^n\ne1\pmod p$ is needed.
Here $n\equiv-18\pmod{p-1}$ and $p>2^{18}-1$, so this condition
holds. Its conversion factor is
\[
 \frac{n2^{n-1}}{2^n-1}
 =\frac{(-19)2^{-19}}{2^{-18}-1}
 =\frac{19}{2^{19}-2}
 =2677869827\pmod p.
 \tag{BC36.2}
\]

For enumeration, choose the primitive root $h=19$.
The complete program proves $p$ prime by trial division, verifies
\[
 p-1=2^3\cdot3^2\cdot5^2\cdot167\cdot27817,
\]
proves the displayed prime factors by trial division, and checks
$19^{(p-1)/q}\ne1\pmod p$ for all their distinct prime divisors
$q$, together with $19^{p-1}=1$.
The first $(p-1)/2=4180895100$ powers $a_i=19^i\bmod p$
therefore give one representative of each pair $\{a,-a\}$.
Pairing the summands of the multiplier-two formula, as in (BC.4),
gives
\[
 S=\sum_{i=0}^{4180895099}
       \left(2\left\lfloor\frac{2a_i}{p}\right\rfloor-1\right)w_i,
 \quad
 w_i=a_i^{-19},\quad
 w_{i+1}=7846251853\,w_i\pmod p.
 \tag{BC36.3}
\]
The roles of the multiplier $2$ and enumerator $19$ are distinct.
The former supplies the signed coefficient and factor (BC36.2);
the latter supplies the exhaustive orbit and weight recurrence.

The exact standard C++17 program
\texttt{higher2/bernoulli\_8361790201.cpp}
implements (BC36.3) in constant memory.
It uses the same proved Montgomery reduction as
Proposition~\ref{prop:bernoulli-residue};
here its bound is again $p^2+2^{64}p<2^{100}<2^{128}$.
The changed orbit multiplication satisfies $19a_i<19p<2^{38}$,
so ordinary unsigned 64-bit multiplication followed by remainder
is exact. Every signed sum update is below $2p$.
All counts, residues, exponents, and comparisons are integers;
floating-point arithmetic appears only in progress timing.
The loop invariant is exactly the residue, weight, and partial sum
in (BC36.3).

The complete run, with all $4180895100$ terms and endpoint checks
$a_{(p-1)/2}=w_{(p-1)/2}=-1$, gives
\[
                     S=7432636175,\qquad
                     \frac{19S}{2^{19}-2}=1827838999\pmod p.
\]
This proves (BC36.1). The source and complete JSON output are included
alongside the full-model certificates; the run took approximately
$31.21$ seconds.

An independently written second program uses multiplier and primitive
enumerator $g=22$, with ordinary unsigned 128-bit modular division.
It independently verifies primality and primitive order, sums all
$4180895100$ terms with the general paired coefficient
$2\lfloor22a/p\rfloor-21$, and checks both endpoints $-1$.
Its complete output is
\[
 S_{22}=1660697019,\qquad
 \frac{n22^{n-1}}{22^n-1}=7951289779,\qquad
 B_n=1827838999\pmod p.
\]
The orbit product is below $22p<2^{38}$, each coefficient times a
weight is below $21p<2^{38}$, and each weight product is below
$p^2<2^{66}$, within its unsigned 128-bit type.
This independent run took approximately $51.28$ seconds.
Its full source and report are
\texttt{higher2/bernoulli\_8361790201\_archive\_g22.cpp/.json}.
It agrees with (BC36.1) using a different multiplier, different
enumeration, and different exact reduction algorithm.
\end{proof}
